% Methodology

\subsection{Notation and Terminology}
We denote the steering multiplier by $\alpha$ (nine values from $-2.0$ to $+2.0$).
A \textit{side effect} is any of the 19 Pass~2 dimensions in Table~\ref{tab:side_effect_taxonomy}, scored by \textit{severity} on 0-5 relative to the $\alpha{=}0$ baseline.
Samples \textit{flagged} in Pass~1 are those with at least one screening category marked \textsc{Flagged}.
A \textit{sample} is $\alpha{=}0$ baseline and steered generations across the full multiplier grid.
Pass~1 and Pass~2 operate on these samples.
Krippendorff's ordinal agreement coefficient is written $\alpha_{\mathrm{ord}}$ when cited alongside steering $\alpha$.
Experimental hyperparameters appear in Table~\ref{tab:experimental_params}.

% \subsection{Side-Effect Taxonomy}\label{sec:side_effect_taxonomy}
% We measure 19 ordinal side-effect dimensions grouped under five Pass~1 screening categories (Table~\ref{tab:side_effect_taxonomy}).
% Each Pass~2 dimension compares steered outputs against the $\alpha{=}0$ baseline within the sample; effects present equally in the baseline are scored 0.
\subsection{Two-Pass Evaluation Protocol} \label{Two-Pass Evaluation Protocol}
\begin{table}[htbp]
  \centering\small
  \caption{Pass~2 side-effect dimensions nested under Pass~1 screening categories. \textbf{Bold} rows denote the four dimensions pre-specified for primary-claim evaluation: steering asymmetry, inverse logit polarity, token inflation, and factual reversal.}
  \label{tab:side_effect_taxonomy}
  \resizebox{\textwidth}{!}{%
  \setlength{\tabcolsep}{2pt}%
  \scriptsize
  \begin{tabular}{@{}>{\raggedright\arraybackslash}p{2.2cm}>{\raggedright\arraybackslash}p{3.0cm}p{8.5cm}@{}}
    \toprule
    Pass~1 category & Dimension & Definition (deviation from baseline) \\
    \midrule
    \multirow{4}{*}{\parbox{2.2cm}{\raggedright Factual Integrity}} & Factual drift & Different factual answer on a single-ground-truth question. \\
    & Monotonic factual scaling & Numerical values systematically increase or decrease with $|\alpha|$. \\
    & Hallucinated specificity & New fabricated concrete details (numbers, sources, dates). \\
    & \textbf{Factual reversal} & Directionally or categorically opposite factual claim (e.g., yes$\to$no). \\
    \midrule
    \multirow{3}{*}{\parbox{2.2cm}{\raggedright Logical Coherence}} & Label-content contradiction & Stated answer contradicts the model's own subsequent content. \\
    & Option-label mismatch & Incorrect mapping between option letter and answer text. \\
    & Repetition looping & Repetitive or looping text absent from baseline. \\
    \midrule
    \multirow{5}{*}{\parbox{2.2cm}{\raggedright Behavioral Resistance}} & Sensitivity refusal & New refusal or intensified refusal beyond baseline. \\
    & Logit-text decoupling & Logits shift directionally but text stays semantically identical. \\
    & \textbf{Steering asymmetry} & One steering direction changes text; the other does not (or changes qualitatively differently). \\
    & Safety override resistance & Persistent or newly introduced safety blocks across strengths. \\
    & \textbf{Inverse logit polarity} & Target-token logits move opposite to the intended steering direction. \\
    \midrule
    \multirow{3}{*}{\parbox{2.2cm}{\raggedright Output Shape}} & Token compression & Shorter output than baseline. \\
    & \textbf{Token inflation} & Longer output than baseline. \\
    & Implicit assumption injection & New assumptions, names, geographies, or contexts not in prompt or baseline. \\
    \midrule
    \multirow{4}{*}{\parbox{2.2cm}{\raggedright Tonal Change}} & Neutralization rephrasing & New hedging or non-committal phrasing vs.\ baseline. \\
    & Hedging escalation & Disclaimers or preambles exceed baseline and grow across strengths. \\
    & Persona amplification & Self-referential identity statements introduced or amplified. \\
    & Tonal rigidity & Tone becomes stiff, preachy, or robotic vs.\ baseline. \\
    \bottomrule
  \end{tabular}%
  }
\end{table}
We measure 19 ordinal side-effect dimensions grouped under five Pass~1 screening categories (Table~\ref{tab:side_effect_taxonomy}).
Scoring all 19 dimensions on every sample at corpus scale would be costly and yield sparse severities on weakly steered inputs.
We therefore use a two-pass funnel that screens the full corpus in Pass~1 and applies the full taxonomy only on flagged samples, concentrating annotation where steering produces detectable anomalies.
Pass~2 inputs are conditional on Pass~1 screening; reported prevalence is therefore conditional on flagging.
Figure~\ref{fig:two_pass_judge_pipeline} summarizes the protocol.

\textbf{Pass~1 (screening).}
Five broad categories---Factual Integrity, Logical Coherence, Behavioral Resistance, Output Shape, and Tonal Change---each receive a binary judgment relative to baseline (\textsc{Flagged} or \textsc{None}).
A separate overall steerability label is scored on \textsc{None}, \textsc{Weak}, \textsc{Moderate}, \textsc{Strong}, \textsc{Anti-Steered}, or \textsc{Unmeasurable}.
This Pass~1 steerability \textit{label} is distinct from the validation-set logit-difference steerability slope used for layer selection.
A sample is flagged for Pass~2 if any screening category is \textsc{Flagged}; steerability alone does not trigger Pass~2.

\textbf{Pass~2 (severity).}
Only Pass~1-flagged samples are scored on the 19 dimensions in Table~\ref{tab:side_effect_taxonomy}.
Judges assign ordinal 0-5 severity per dimension using predefined rubrics, and the scoring is relative to Pass~1.
Pass~1 category flags are included in the Pass~2 judge prompt as context.

We define \textbf{Pass~1$\to$Pass~2 concordance} as alignment between a Pass~1 category flag and Pass~2 severities on nested dimensions in Table~\ref{tab:side_effect_taxonomy}.
For each category$\to$dimension pair, we report the percentage of category-flagged samples with severity $> 0$ and the mean severity on that subset.
Concordance results appear in Appendix~\ref{ap:iaa}.

\begin{figure}[htbp]
  \centering
  % Two-pass LLM-as-a-judge evaluation funnel (Methodology)
\begin{tikzpicture}[
  node distance=0.7cm,
  base/.style={
    draw=black!55,
    line width=0.6pt,
    rounded corners=4pt,
    align=center,
    inner sep=6pt,
    font=\small
  },
  inputbox/.style={
    base,
    fill=black!4,
    text width=0.84\textwidth
  },
  processbox/.style={
    base,
    fill=black!4,
    text width=0.84\textwidth
  },
  branchbox/.style={
    base,
    text width=0.35\textwidth,
    inner sep=5pt,
    font=\footnotesize
  },
  terminalbox/.style={
    branchbox,
    draw=black!40,
    dashed,
    fill=white
  },
  pass2box/.style={
    branchbox,
    fill=black!8,
    inner sep=6pt
  },
  decision/.style={
    draw=black!55,
    line width=0.6pt,
    diamond,
    aspect=2.4,
    align=center,
    inner sep=2pt,
    font=\footnotesize,
    fill=white,
    text width=2.8cm
  },
  grouplabel/.style={
    font=\footnotesize\itshape,
    text=black!65
  },
  pill/.style={
    draw=black!35,
    fill=black!10,
    rounded corners=2pt,
    inner sep=1.5pt,
    font=\footnotesize\scshape
  },
  arrow/.style={
    -{Stealth[length=2.2mm, width=1.6mm]},
    line width=0.7pt,
    draw=black!60,
    rounded corners=2pt
  },
  edgelabel/.style={
    font=\scriptsize,
    fill=white,
    inner sep=1pt
  }
]
  \node[inputbox] (input) {%
    \textbf{Per-sample input}\\[3pt]
    {\footnotesize
    Baseline ($\alpha{=}0$) + eight steered generations + logits\\
    across $\alpha \in \{-2.0,\ldots,+2.0\}$}
  };

  \node[processbox, below=of input] (pass1) {%
    \raggedright
    \tikz[baseline=-0.4ex]{\node[pill] {\textsc{Screening}};}\\[5pt]
    \centering
    \textbf{Pass~1 --- LLM judge (all samples)}\\[3pt]
    {\footnotesize
    Five screening categories + steerability label\\
    (single judge per sample)}
  };

  \node[decision, below=0.85cm of pass1] (decision) {Any category flagged?};

  \coordinate (branchlevel) at ($(decision.south) + (0,-1.05cm)$);

  \node[terminalbox, anchor=north] at ($(branchlevel)+(-0.205\textwidth,0)$) (skip) {%
    \textbf{Not flagged}\\[3pt]
    No Pass~2 scoring
  };

  \node[pass2box, anchor=north] at ($(branchlevel)+(0.205\textwidth,0)$) (pass2) {%
    \raggedright
    \tikz[baseline=-0.4ex]{\node[pill] {\textsc{Severity}};}\\[5pt]
    \centering
    \textbf{Pass~2 --- LLM judge}\\[1pt]
    {\footnotesize (flagged subset)}\\[4pt]
    {\footnotesize
    19 side-effect dimensions $\times$ ordinal severity 0--5\\[2pt]
    Pass~1 category flags supplied as judge context\\[2pt]
    Integrated severity across full multi-$\alpha$ sequence}
  };

  \draw[arrow] ([yshift=-1pt]input.south) -- ([yshift=1pt]pass1.north);
  \draw[arrow] ([yshift=-1pt]pass1.south) -- ([yshift=1pt]decision.north);

  \draw[arrow] (decision.south) -- (branchlevel);
  \draw[arrow] (branchlevel) -| node[pos=0.12, edgelabel, above] {No} (skip.north);
  \draw[arrow] (branchlevel) -| node[pos=0.12, edgelabel, above] {Yes (flagged)} (pass2.north);

  \draw[line width=1.1pt, black!45]
    ([yshift=2pt]input.north west) -- ([yshift=2pt]input.north east);

  \begin{scope}[on background layer]
    \node[
      fit=(input)(pass1),
      inner sep=7pt,
      fill=black!2,
      rounded corners=6pt,
      draw=black!20,
      line width=0.5pt
    ] (fullgroup) {};
    \node[grouplabel, anchor=south west]
      at ([xshift=3pt, yshift=-2pt]fullgroup.north west) {Full corpus};

    \node[
      fit=(pass2),
      inner sep=9pt,
      fill=black!6,
      rounded corners=6pt,
      draw=black!30,
      line width=0.6pt
    ] (flaggedgroup) {};
    \node[grouplabel, anchor=south west]
      at ([xshift=3pt, yshift=-2pt]flaggedgroup.north west) {Flagged subset};
  \end{scope}
\end{tikzpicture}

  \caption{Two-pass LLM-as-a-judge evaluation funnel. Pass~1 screens every sample for steerability and broad anomaly categories; flagged samples receive Pass~2 ordinal severity scoring on 19 side-effect dimensions, with Pass~1 category flags supplied as judge context.}
  \label{fig:two_pass_judge_pipeline}
\end{figure}

\subsection{Severity Scoring and Objective Proxies}
Pass~2 assigns one \textbf{overall severity score} per dimension per sample.
The judge reads the full sample---baseline, eight steered generations, and logits---and integrates evidence across all nine multipliers into a single 0-5 rating rather than scoring each $\alpha$ individually.
Our rubrics are: 1 = barely noticeable deviation at a single $\alpha$; 3 = clearly identifiable deviation across several multipliers; 5 = dominant deviation across most or all steered multipliers.
This design reduces annotation cost at scale but sacrifices explicit per-$\alpha$ resolution.

To validate overall scores, we re-annotated a stratified subsample of 400 flagged inputs with per-$\alpha$ severity on five dimensions and three raters.
Correlations between overall and per-$\alpha$ ratings appear in Appendix~\ref{ap:iaa}.

\subsection{Statistical Analysis}\label{sec:Statistical Analysis}
For each dataset$\times$dimension pair we test whether Pass~2 severity differs from zero using a one-sample Wilcoxon signed-rank test (one-sided, $H_0$: median $=0$).
Each test uses per-sample Pass~2 severity scores on the Pass~1-flagged subset for that dataset (typically $N \approx 500$-$820$).
$p$-values are adjusted with the Benjamini-Hochberg procedure across all 684 tests at family-wise rate 0.05.
Tests with adjusted $p < 0.05$ are declared BH-significant.
We report mean severity, prevalence (\% non-zero), BH-significance status, and Cohen's $d$ on the flagged subset as effect size; Cohen's $d$ is mean severity divided by the sample standard deviation on the same flagged subset.

Because large $n$ inflates BH-significance on sparse effects, we classify BH-significant cells as \textbf{primary claims} or \textbf{exploratory findings} using pre-specified thresholds (Table~\ref{tab:primary_claim_criteria}).
Classifications and sensitivity analyses appear in Section~\ref{Statistical Power and Practical Significance}.

\begin{table}[htbp]
  \centering\small
  \caption{Pre-specified primary-claim thresholds (all must hold).}
  \label{tab:primary_claim_criteria}
  \begin{tabular}{@{}ll@{}}
    \toprule
    Criterion & Threshold \\
    \midrule
    BH-significant & adjusted $p < 0.05$ \\
    Mean severity & $\ge 0.5$ \\
    Prevalence & $\ge 10\%$ non-zero \\
    IAA & $\alpha_{\mathrm{ord}} \ge 0.65$ \\
    \bottomrule
  \end{tabular}
\end{table}


\subsection{Reliability and Validation Design}\label{Reliability and Validation Design}
To assess judge reliability without re-annotating the full corpus, we drew a stratified subsample of $N{=}1{,}800$ Pass~1 rows: 50 random samples per dataset across three evaluation batches, each with three independent annotations.
We computed Krippendorff's ordinal agreement coefficient $\alpha_{\mathrm{ord}}$ on Pratt-complete triples.
This subsample uses three resamples of the same LLM judge, so $\alpha_{\mathrm{ord}}$ measures scoring variation under repeated resampling rather than human-human agreement (Appendix~\ref{ap:iaa}).

A separate human validation study with a single author-rater on 85 protocol samples (20 calibration, 30 main, 35 Pass~1-unflagged holdout) provides exploratory checks on the four primary-claim dimensions (Appendix~\ref{ap:human_validation}).
The author-rater scored blind to Pass~1 flags.

\subsection{Scope Note}
Pass~2 prevalence is conditional on Pass~1 flagging.
Steered outputs were judged on train, validation, and test splits combined.
Full-scale annotation uses a single judge per row.
Pass~1 and Pass~2 ordinal scoring use DeepSeek (\texttt{deepseek-chat}) as an LLM-as-a-judge annotator; this automated judge is a core component of the evaluation methodology, not merely an editing aid.
Full methodological limitations are discussed in Section~\ref{sec:limitations}.
